\documentclass[12pt]{article}
\usepackage{../../../assignment}
\usepackage{../../../math}
\usepackage{../../../uark_colors}
\hypersetup{
  colorlinks = true,
  allcolors = ozark_mountains,
  breaklinks = true,
  bookmarksopen = true
}

\begin{document}
\begin{center}
  {\Huge\bf Final - Fall 2025}

  \smallskip
  {\large\it ECON 5783 — University of Arkansas}
\end{center}
\medskip

\begin{enumerate}
  \item A researcher wants to evaluate whether a school having an arts program benefits children's educational outcomes. 
  The researcher surveys schools and records their average GPA ($Y_i$) and whether or not they have an arts program ($D_i$).
  \begin{enumerate}
    \item In this context, describe in words what the potential outcomes, $Y_i(0)$ and $Y_i(1)$, are.
    
    \medskip
    \item Do you think the difference-in-means estimator would be biased in this context? 
    If so, which direction do you the treatment effect estimate is biased?
    
    \medskip
    \item Why might a table of summary statistics consisting of sample means of observables, broken down by $D_i = 0$ and $D_i = 1$, be useful in this setting?
  \end{enumerate}
\end{enumerate}

\bigskip\bigskip
%% https://www.christopher-neilson.com/work/documents/UPK/w33038.pdf
\noindent The following questions are about Rosenthal and Strange (2008, Journal of Urban Economics). 
They want to understand if having a higher density of firms in an area increases the average wages of workers. 

\begin{enumerate}
  \setcounter{enumi}{1}
  \item Consider regressing the average wages of a worker in a city on the density of firms in the area. 
  Say you find a large, positive estimate.
  Using the concept of omitted variables bias, give an example of why we should not believe we identified the causal effect of density.

  \medskip
  \item Rosenthal and Strange use as an instrument the proportion of soil in the city that contains sedimentary rock. 
  The idea being that sedimentary soil is more difficult to build on which causes less density. 
  Argue why the exclusion restriction might be reasonable in this context.

  \medskip
  \item The first-stage regression of density on the sedimentary rock share has a $F$-statistic of 1.907. 
  Describe why this might cause a problem when conducting inference on our 2SLS estimate.
\end{enumerate}



\newpage
\noindent The following questions are about Leah Boustan's work on court-enforced school desegregation (2012, AEJ Applied).
They ask whether forced school desegregation in the city-center led to a larger decline in the housing price in the city-center.
The unit of observation is a metro area and they are observed in 1960, 1970, and 1980 (post-period).
Below is the primary DID estimates.

\bigskip
\begin{enumerate}
  \setcounter{enumi}{4}
  \item The row with $\Delta$ 1970--1980 presents the difference-in-differences estimate. 
  What do we have to assume about the cities that were and were not placed under court orders in order to belive the 5.8\% decrease in home prices estimate?
  
  \medskip
  \item Why is the row labeled ``$\Delta$ 1960--1970'' helpful for justifying the causal assumption?
  
  \medskip
  \item Boustan suggests that cities that had larger shares of Black residents might have had larger suburbanization rates during this time (`White flight'). 
  Moreover, cities that put under court-ordered forced desegregation had larger declines in city housing prices.
  Why might this represent concerns for our parallel trends assumption?

  \medskip
  \item Say you are willing to assume parallel trends \emph{conditional on} the share of Black residents in 1970, i.e. $\Delta Y_i(0) \Perp D_i \ \vert \ X_i = x$.
  Describe how you would estimate a treatment effect using the regression adjustment strategy. 
\end{enumerate}

\bigskip
\begin{table}[hb]
\centering
\caption{School Desegregation and Relative City Housing Prices at the District Border, 1960--1980}
\label{tab:desegregation}
\begin{tabular}{lccc}
\toprule
Dependent variable & Placed under & Not placed under & \\
$= \ln(\text{housing value})$; & court order & court order & \\
coefficient on & during 1970s & during 1970s & Difference \\
\midrule
% 1970 & $-0.047***$ & $-0.026*$ & $-0.021$ \\
%      & (0.014) & (0.015) & (0.020) \\[0.5em]
% 1980 & $-0.097***$ & $-0.023$ & $-0.073**$ \\
%      & (0.028) & (0.022) & (0.035) \\[0.5em]
$\Delta$ 1970--1980 & $-0.065**$ & $-0.007$ & $-0.058**$ \\
     & (0.024) & (0.015) & (0.028) \\[0.5em]
\textit{Pre-trend:} & & & \\
$\Delta$ 1960--1970 & $-0.023*$ & $-0.022$ & $-0.001$ \\
     & (0.013) & (0.017) & (0.022) \\
\bottomrule
\end{tabular}
\end{table}

\end{document}
